Elke fetish of seksuele obsessie is eigenlijk een gefantaseerde oplossing voor een angst in de echte wereld. Emma houdt van uniformen want ze is bang voor mensen in uniforms omdat ze haar slecht hebben behandeld, terwijl ze door zelf te spelen in sex de rollen omdraait en het juist leidt tot een goede behandeling. Het onbereikbare koude autoriteit is een bereikbare warme geworden die ons niet intimideert of belemmert maar geruststelt en helpt. Sex geeft hints naar een utopie waar kracht, organisatie, netheid en structuur er is om ons meer op ons gemak te laten voelen.

Emmy houdt van onderdanigheid omdat ze al jong leerde om onafhankelijk te zijn. Ze groeide op in een individualistisch gezin die afhankelijkheid van de landheer als onwenselijk ervaarde en daarom haar duwde naar een solitaire volwassenheid van Emmy. Daarom verlangd ze naar het wegnemen van haar verantwoordelijkheid en keuzevrijheid, zich over te geven aan iemand die de regie neemt, de druk van het dagelijks leven op zich neemt en de last op zich neemt haar ook helpt drempels op sexueel of ander vlak te overwinnen. Later in het boek vindt ze die persoon in de oude wijze. Daarvoor leert ze dat de jager daarin niet de juiste persoon is maar dat ze iemand nodig heeft die goed en aardig is, hoewel ze uit zelfbescherming en angst dit beeld toch maar op de jager plakt. Ze wil eigenlijk leren om lijden te verdragen in plaats van weg te dragen. Om sterk te blijven terwijl je zelf wil dat iemand je een slet of snol noemt

Ze beleefde op een vreemde manier
toch een bepaalde vorm van plezier.

Alsof ze met hem meegenoot
toen hij zijn dna in haar spoot

en zich voldaan uit haar rolt.
Alsof ze zich met hem versmolt.

Ze samen met de jager Emmy nam
en samen met de jager klaarkwam.

In haar verdere reflectie
op deze bijzonder connectie

dacht ze aan de pesters terug
tegen de muur met haar rug.

Hoe ze ook daar kon meeleven.
Zich zodoende aan hen kon geven.

Haar eigen belang niet meer telde.
De vernedering haar niet meer kwelde

Emmy zich zelfs kon voorstellen
hoe de opwinding kon aanzwellen

als zij zo hulpeloos voor hen stond
van angst genageld aan de grond.

Hoe lekker het vernederen kan zijn
door te ejaculeren in haar bijzijn.

Hoezeer machtsmisbruik aan je kan trekken.
Ze leerde haar onderdanigheid ontdekken.

Je zelf uiteindelijk willen opofferen.
Door je lichaam te slachtofferen.

Geen onderscheid meer willen ervaren.
Het gevoel van onveiligheid willen bedaren.

Niet een ree, maar een roedel zijn.
Niets krijgt je dan meer klein.

Geen mongool met zijn zieligdoenerij.
Geen Herder met zijn vleierij.

Emmy was ze allemaal tegelijk.
Werd de kille wereld een hemelrijk.

Als ze haar ego tijdelijk verloor,
brak een diepe vrede in haar door.

Kon ze een penibele situatie willen
zoals het prikken tussen haar billen.

Was ze net zo kalm als de jager.
Bevredigd als slachtoffer en belager.

Door Emmy vol aan hem te geven
kon ze het genot met hem beleven.

De toestand van totale eenheid
gaf haar kalmte en tevredenheid.

Deze akelige situatie
was haar eigen creatie.

Alsof het haar woorden waren,
hoorde ze de jager verklaren,

hoe hij had genoten van dit spel.
Ze wachtte op het woord “vaarwel”.

Deze houding was gemakzuchtig.
Haar manier bijzonder luchtig.

De jager gedroeg zich ongepast.
Maar zij droeg zijn schuldenlast.

Het was niet aleen bangheid.
Ze negeerde haar boosheid.

De jager was over de grens gegaan.
Dat inzicht kon zij niet toestaan.

De jager kneep in haar billen.
Ze wilde het uitgillen.

Een film in haar koppie
greep echter haar attentie.

Ze was een princes
en kreeg dansles.

Mannen stonden eromheen
en betasten soms een been

of knepen in een bil
terwijl ze dat niet wil.

Ze herinnerde dit vorig leven
alsof ze het opnieuw kon beleven.

Ze leerde zich slap te houden
als mannen zich niet inhouden.

Eentje daarvan wacht haar op
ging naar de kleedruimte voorop.

Hij greep haar handen
die zich flink inspanden

om zich los te rukken.
Dat wilde niet lukken.

Hij duwde haar op de grond
en kust haar op haar mond.

Ze worstelt hevig.
Zijn grip is te stevig.

Hij trekt haar jurk omhoog.
Hij houdt zijn kruit droog.

Ze probeert te schoppen.
Hij is niet te stoppen.

Dan komt zijn hoogtepunt.
Rust is hem niet gegund.

Ze dreigt met openbaarmaking.
Hij snijdt in één beweging

haar keel met mes door
waarna zij haar leven verloor.

Ze zag een fel licht.
Haar ogen gingen dicht.

Ze stapte kalm vooruit
richting het stemgeluid

dat haar verder leidde
en met woorden voorbereidde

op wat er komen ging.
Ze ervaarde een bevalling.

Ze hoorde zichzelf huilen.
Kon zich niet verschuilen

in de moederschoot.
Het was blijk en bloot

dat ze was herboren.
Als Emmy was geboren.

Toen kwam het inzicht
waarom ze makkelijk zwicht

voor de wensen van de man.
Omdat ze niet anders kan.

Ze had eenmaal geaccepteerd,
uit een vorig leven geleerd,

dat verweer het erger maakt.
Ze daarom elk verzet staakt.

Gooide herinneringen overhoop

de stellingen van Tannen.
Die stonden op gespannen

voetEmmy’s intuitie.
Volgens de benedictie

van de abt wees voorgevoel
naar ons werkelijk levensdoel.

De abt was intelligent
en een grote sterke vent

die het mannenklooster leidde.
Ook zijn boodschap verspreidde

in het klooster van de nonnen.
Hij was daar ooit mee begonnen

toen de Abdis werd verweten
dat ze het flink liet afweten.

Ze wist steeds verder van huis te raken
in haar poging om verlangens te staken.

Volgens de abt was ook dit een verlangen
waardoor ze in haar missie bleef hangen.

Fellatio, aftrekken en bruinwerken
Ze kon niet anders dan opmerken

dat dit geheel niet strookt
met waarmee ze was bestookt

in het klooster en de bijbel.
Het had zeker geleid tot heibel

als ze dit de abt had gemeld.
Die had haar vast verteld

dat Tannen haar heeft bedrogen.
Haar vals heeft voorgelogen

om anderen een dienst te bewijzen.
Het vulde Emmy met afgrijzen.

Ze was geniepig onteerd.
Gelukkig had ze ook geleerd;

“Zalig zijn de armen van geest”.
Was zij niet gewoon trouw geweest?

Kwam niet haar decorumverlies
door een misleidend advies

van iemand met religieuze autoriteit.
Die haar schaamteloos had misleid?

Was ze niet nederig over haar niet willen
om zichzelf niet over het paard te tillen?

Uit een enquête onder meer dan 4.000 Amerikanen, besproken in Psychology Today, blijkt dat 69,8 % van de Amerikanen seksuele fantasieën heeft. De voornaamste reden hiervoor is nieuwsgierigheid naar verschillende sensaties en seksuele ervaringen. Naast het opwekken van seksuele opwinding, wat de meest genoemde reden was, gaven deelnemers aan dat ze fantaseren om onvervulde seksuele behoeften te bevredigen, even aan de realiteit te ontsnappen, of om een sociaal taboe seksueel verlangen te uiten en te vervullen.

Maar wat maakt een seksueel verlangen precies een taboe? Het zal je misschien verbazen dat veel van wat de maatschappij als taboe beschouwt, eigenlijk heel gewone seksuele fantasieën zijn. Hier zijn enkele voorbeelden van schijnbaar taboe fantasieën die verrassend vaak voorkomen.

Dominantie en Onderwerping

De fantasie om gedomineerd te worden of zelf te domineren komt vaker voor dan je zou denken. Sociaal psycholoog Justin Lehmiller ondervroeg 4.000 Amerikanen en ontdekte dat 93 % van de vrouwen en 81% van de mannen fantaseerde over gedomineerd worden. Andersom fantaseerde 85% van de mannen en 76% van de vrouwen over seksueel dominant zijn.

Misschien heb je wel eens gehoord van BDSM, wat staat voor bondage, discipline, dominantie, submissie (onderwerping), sadisme en masochisme. Bij BDSM is vaak een psychologische component betrokken. Veel mensen houden bijvoorbeeld van psychologisch rollenspel waarin een machtsuitwisseling centraal staat. De invulling van dit soort fantasieën kan sterk variëren. Een artikel in Glamour legt uit dat het iets subtiels kan zijn, zoals een vertrouwde partner die je polsen vasthoudt terwijl hij of zij je kust, tot extreme BDSM met pijn, vernedering of wat je kinky hart ook maar begeert.

Een studie gepubliceerd in de Journal of Sexual Medicine toonde aan dat 65 % van de mensen die zich als vrouw identificeren, verlangt om gedomineerd te worden, terwijl 46% van hen toegeeft wel eens gefantaseerd te hebben over het domineren van iemand anders. The Washington Post plaatste echter een kanttekening bij deze studie: de deelnemers aan seksstudies zijn vaak seksueel ervarener dan gemiddeld. De auteurs van de studie gaven zelf ook aan dat de steekproef niet representatief was voor de algemene bevolking. Toch doet de populariteit van boeken als "Fifty Shades of Grey" je afvragen hoe wijdverspreid deze fantasieën werkelijk zijn.

Volgens auteur en gecertificeerd sekscoach Gigi Engle kan het idee van seksuele onderdanigheid opwindend zijn voor mensen die buiten de slaapkamer altijd de controle hebben. Omgekeerd kan het idee van de controle hebben juist opwindend zijn vanwege het taboe-karakter van ruwe seks en het gevoel van autoriteit.

Bondage

Bondage, een onderdeel van BDSM, wordt ook vaak als taboe gezien, maar blijkt een veelvoorkomende seksuele fantasie te zijn. Opvoeder en seksuoloog Midori legt uit dat bondage het gevoel van gegeven en genomen macht in dominantie- en onderwerping spel kan versterken. Ze stelt dat fysiek vastgebonden worden met touw, handboeien of zelfs een stropdas, iemand met deze fantasie een totale lichamelijke en mentale bevrijding van verantwoordelijkheid kan geven. Duidelijke toestemming en leren hoe je deze fantasie veilig kunt uitleven, zijn absoluut noodzakelijk wanneer je dit verlangen met een partner in de praktijk brengt.

Exhibitionisme

Sommige mensen fantaseren over strippen voor een publiek, seks hebben voor de camera, of masturberen in het bijzijn van iemand anders, een groep of hun partner. Deze fantasieën kunnen gaan over het pronken met jezelf of de spanning van het betrapt worden. AASECT-gecertificeerd sekstherapeut Indigo Stray Conga definieert exhibitionisme als het ervaren van seksuele opwinding door de fantasie (of de daad) van naakt of tijdens seksuele handelingen geobserveerd te worden.

Stray Conga benadrukt dat gezond exhibitionisme een sekspositieve viering van het erotische is en niet verward moet worden met de exhibitionistische stoornis. Ze merkt ook op dat het een extreem veelvoorkomende fantasie is. NYU-professor en seksexpert Zhana Vrangalova beschrijft de bevindingen van een studie uit 2015, gepubliceerd in de Journal of Sexual Medicine. Deze studie toonde aan dat 66 % van de mannen en 50% van de vrouwen fantaseert over seks op een openbare, publieke plaats. Daarnaast fantaseert 82% van zowel mannen als vrouwen over seks op een ongebruikelijke plek, zoals een kantoor of een openbaar toilet. "De meesten van ons hebben wel een klein exhibitionistisch trekje," legt ze uit. "Dit bestaat, zoals alles in de psychologie, op een spectrum en is volkomen normaal en gezond."

Groepsseks

Een seksuele fantasie waarbij meer dan één andere partner betrokken is, is ook vrij normaal. Een studie gepubliceerd in het tijdschrift Personality and Individual Differences ondervroeg 788 Britse volwassenen in vier leeftijdsgroepen. Ze ontdekten dat mannen vaker fantaseerden over seks met meer dan één persoon, vaak met anonieme partners. Bij vrouwen stelde de studie dat zij, als aandeel van hun totale fantasieën, vaker fantaseerden over partners van hetzelfde geslacht en over beroemde partners.

Of het nu gaat om een trio of een orgie, fantasieën met meerdere partners lijken heel gewoon te zijn. Een studie uit 2017, gepubliceerd in Archives of Sexual Behavior en besproken in een artikel van Insider, toonde aan dat van de 274 Canadese universiteitsstudenten 64 % interesse had in seks met meerdere partners.

Het feit dat deze seksuele fantasieën, die vaak als taboe worden gezien, zo veel voorkomen, kan voor sommige mensen een opluchting zijn. Gecertificeerd sekscoach Gigi Engle legt uit dat deze erotische gedachten vaak volkomen normaal zijn. Volgens Engle geldt: hoe meer we praten over seksuele fantasieën en het gesprek erover normaliseren, hoe minder we onszelf zullen veroordelen voor het hebben van spannende, erotische gedachten.

Oké, hier is de Nederlandse vertaling van het YouTube-transcript, omgezet in een lopend verhaal zonder de tijdscodes.

Inleiding
Voor sommigen is een blinddoek en een lichte 'spanking' al super kinky, terwijl anderen het slechts als voorspel zien en pas echt opgewonden raken wanneer de 'ball gag' en leren peddels tevoorschijn komen. Maar wat als seksuele fetisjen een duistere wending nemen?
Voor sommige mensen kan een fetisj veranderen in een obsessie die bekend staat als parafilie of een parafilische stoornis. Degenen met een parafilische stoornis hebben vaak intense, aanhoudende seksuele fantasieën die te maken hebben met de psychologische nood, het letsel of zelfs de dood van een andere persoon. In sommige gevallen komen deze seksuele afwijkingen vaker voor bij mannen en omvatten ze gedragingen zoals exhibitionisme, voyeurisme en frotteurisme. Hoewel fetisjisme onmiskenbaar een duistere kant heeft, raken sommige individuen opgewonden door veel meer dan zwepen en kettingen. Laten we het hebben over wat die zijn.
Hybristofilie
Stel, je neemt iemand mee op date en je hoopt op wat fysiek contact. Een goed idee is om naar een pretpark of een enge film te gaan. Dit zorgt ervoor dat de adrenaline en endorfine door hun lichaam gieren, waardoor ze waarschijnlijk dicht bij je in de buurt blijven uit angst. Dit is een voorbeeld van een gezonde manier om de chemie van je lichaam te gebruiken om gevoelens te versterken.
Dan zijn er mensen met hybristofilie. Deze individuen hebben hun verlangen om rebels te zijn tot het uiterste doorgevoerd. Ze voelen zich seksueel aangetrokken tot criminelen, criminele activiteiten of iets wat ook maar enigszins illegaal is. Een van de bekendste voorbeelden tot op de dag van vandaag zijn de vrouwen die zich aangetrokken voelden tot Ted Bundy. Naar verluidt ontving hij honderden liefdesbrieven van vrouwen terwijl hij in de gevangenis zat, kreeg hij een huwelijksaanzoek, trouwde hij met een van hen en zou hij de vader van haar kind zijn.
Gevallen van hybristofilie zijn het klassieke "ik kan hem/haar wel veranderen"-verhaal, maar dan in een gevaarlijke vorm. Denk aan de Joker en Harley Quinn. De een raakt aangetrokken tot de ander door de pure adrenaline en het nieuwe. Zodra ze elkaar beter leren kennen, komt het verlangen om de ander te 'repareren' op en raakt dit verweven met hun gevoelens. En dat is waar het gevaarlijk wordt. Onthoud: iemand die echt om je geeft, zal je nooit opzettelijk in gevaar brengen en zal altijd willen dat je veilig bent.
Frotteurisme
Stel je voor dat je in een bus zit. Elke stoel is bezet en er staan zelfs passagiers in het gangpad. Jij bent een van de staande passagiers. Wanneer er een nieuwe passagier instapt, schuift iedereen op om ruimte te maken. Iemand met frotteurisme zou echter, in plaats van ruimte te maken, opzettelijk en op een ogenschijnlijk nonchalante manier tegen de andere staande passagiers aanbotsen en schuren. Dit is frotteurisme, beter bekend als 'aanranding door betasting' (groping).
Dit is wanneer een individu seksuele opwinding of plezier haalt uit het tegen een andere persoon aanschuren of deze aanraken met de bedoeling opgewonden te raken. Dit gedrag is natuurlijk niet alleen sociaal onaanvaardbaar zonder toestemming van alle partijen, maar wordt ook beschouwd als aanranding. Deze seksuele fetisj kan zich ontwikkelen door jeugdtrauma, specifiek door een gebrek aan fysieke genegenheid van een verzorger of door blootstelling aan seksueel materiaal op jonge leeftijd.
Somnofilie
Herinnert iemand zich die scène in Twilight waarin Edward de kamer van Bella binnensluipt om naar haar te staren terwijl ze slaapt? Iemand gediagnosticeerd met somnofilie raakt opgewonden door seks te hebben met of iemand seksueel aan te raken die slaapt. Tenzij dit in een eerder gesprek is overeengekomen, kan er op die momenten geen toestemming worden gegeven.
Het Journal of Sleep Disorders and Therapy geeft een aantal mogelijke redenen waarom individuen somnofilie ontwikkelen. Dit omvat het gevoel weinig tot geen controle te hebben over hun eigen leven, onopgeloste problemen met dominantie of kwetsbaarheid, of als een reactie op jeugdtrauma. Het observeren van of omgaan met iemand die bewusteloos is, betekent dat die persoon niet kan protesteren of afwijzen, wat de 'jager' een vals gevoel van controle en macht geeft.
Weet alsjeblieft dat het niet oké is als iemand je aanraakt of bekijkt op een seksuele manier terwijl je bewusteloos bent, tenzij alle partijen hiermee hebben ingestemd terwijl ze bij bewustzijn waren. Als jij of iemand die je kent een vergelijkbare ervaring heeft gehad zonder toestemming te geven, bespreek dit dan alsjeblieft met vertrouwde dierbaren, medische of geestelijke gezondheidsprofessionals, of een andere vertrouwde autoriteit.
Necrofilie
Wat als iemands angst voor afwijzing zo intens is dat ze zich aangetrokken voelen tot iets of iemand die niet leeft? Necrofilie, het aangetrokken zijn tot een lijk, komt voort uit het idee van het hebben van een partner of object dat zich niet zal verzetten of hen zal afwijzen.
Havelock Ellis suggereerde in zijn boek Studies in the Psychology of Sex uit 1903 dat necrofilie en algolagnie (seksueel genot uit pijn) met elkaar verbonden zijn. Beide houden in dat negatieve gevoelens zoals angst, vrees en afkeer worden omgezet in opwinding en verlangen. Een dode kan echter geen toestemming geven, en seksueel genot uit een lijk wordt over het algemeen beschouwd als een diepe vorm van respectloosheid.
Dus, welke andere factoren leiden tot een dergelijke aantrekkingskracht? Onderzoek suggereert ook dat sommige mensen zich bezighouden met necrofilie omdat ze proberen opnieuw verbinding te maken met een overleden geliefde of partner. De aantrekkingskracht ligt vaak in het samenzijn met iemand die hen niet zal afwijzen of tegenspreken, wat voor sommigen veiliger voelt. Voor hen kan alleen al het zien van een lijk seksueel opwindend zijn. Het is een manier om eenzaamheid te bestrijden, troost te zoeken en controle te hebben over iemand die zich niet kan verzetten.
Naast deze hoofdmotieven zijn er minder vaak voorkomende redenen, zoals angst voor het andere geslacht, het ervaren van bevelshallucinaties of een verlangen naar volledige controle. Erich Fromm suggereert dat necrofilie kan voortkomen uit een poging van een individu om een gebrek aan zelfidentiteit en authenticiteit te compenseren.
Kortom, als jij of iemand die je kent met dit soort problemen worstelt, is het heel belangrijk om met een professional in de geestelijke gezondheidszorg te praten. Zij kunnen helpen om deze gevoelens te verwerken en gezondere manieren te vinden om met afwijzing, eenzaamheid en controleproblemen om te gaan.
Auto-erotische verstikking
Auto-erotische verstikking is wanneer mensen proberen hun seksuele genot te verhogen door de luchttoevoer naar de hersenen te beperken, meestal door wurging of verstikking. Dit kan leiden tot een euforische 'high' of intense opwinding. Sommige mensen proberen dit tijdens het masturberen met behulp van touwen, koorden of plastic zakken, in de hoop hun ervaring te intensiveren. Dit is echter ongelooflijk gevaarlijk en kan snel leiden tot bewusteloosheid, verstikking of zelfs de dood.
Een beroemd geval is de dood van David Carradine in 2009, bekend van zijn rol in Quentin Tarantino's Kill Bill: Volume 2. Hij werd dood aangetroffen in zijn hotelkamer in Bangkok. Aanvankelijk werd gedacht aan zelfmoord, maar twee pathologen stelden later vast dat het een ongeluk was. Zijn ex-vrouwen onthulden dat hij een geschiedenis had van zelfbondage, een andere riskante fetisj.
Zelfs met een partner wordt seksuele verstikking gezien als een van de meest riskante 'kinks' in de BDSM-gemeenschap. Het brengt ernstige gevaren met zich mee, waaronder het risico op overlijden of blijvend letsel, zelfs als alles correct lijkt te worden uitgevoerd.
We zijn hier niet om iemands fetisjen of voorkeuren te veroordelen. Als je een eerlijk en open gesprek hebt gehad met je romantische partners over je fetisjen, ga dan vooral je gang. Als je echter gelooft dat jij of een dierbare lijdt aan een duistere seksuele fetisj, praat dan alsjeblieft met een medische of geestelijke gezondheidsprofessional.

Een korte opmerking over de inhoud: deze aflevering behandelt enkele behoorlijk zware onderwerpen, waaronder seksueel geweld, kindermisbruik en enkele korte, grove vermeldingen. Als dat niets voor jou is, geen probleem. Sla deze dan gewoon over.
De Gekaderde Kijk op Kink
Als je aan kink denkt, waar denk je dan aan? Als je er voornamelijk mee in aanraking komt via populaire media, dan is je beeld van kink waarschijnlijk verdeeld in twee smaken. Er is de leuke, sexy en soms gekke soort: de sitcom-moeder die haar huwelijk wil opfleuren met pluizige handboeien en een blinddoek, of de comédienne die haar liefde voor 'dirty talk' en een pak op de billen belijdt. En dan is er die ene scène waarin een jonge Chris Evans slagroom op zijn tepels spuit en een banaan in zijn achterste stopt – ja, zoek het maar op.
Daarnaast is er de gevaarlijke soort: de in leer geklede dominatrix die de heteroseksuele mannelijke hoofdpersoon in haar kerker opsluit, wat leidt tot allerlei capriolen. De enge schoenenwinkel-eigenaar met een voetfetisj. De seriemoordenaar die zich graag in vrouwenkleding hult. De manier waarop we leren naar kink te kijken, is ofwel als een acceptabele toevoeging aan een "normaal", heteroseksueel, monogaam en privé seksleven, of als iets pathologisch, iets wat wordt gedaan door afwijkende mensen met duistere bedoelingen.
Antropologe Margot Weiss stelt dat dit ons tot toeschouwers van kink maakt, die vanuit een veilige, normale en bevoorrechte positie labelen en oordelen, in plaats van de genuanceerde realiteit ervan te begrijpen. We houden kink op een afstand. In deze miniserie wil ik die afstand verkleinen. Ik ga je laten zien dat kinks en fetisjen veel meer zijn dan de tweedeling tussen vrijgezellenfeest-cadeaus en zedenregisters. Het is een spectrum, en jij bevindt je er ook op.
Ik noem deze serie "Philias," naar het Griekse woord en wetenschappelijke achtervoegsel voor liefde. Dat achtervoegsel komt meestal aan het einde van het woord parafilie, wat per definitie een abnormaal verlangen is, een afwijkende vorm van liefde. Maar zoals we zullen leren, is de wetenschap er notoir slecht in om te definiëren wat normaal is. In plaats daarvan zijn de vele kinks en fetisjen die in de wereld bestaan een prachtige illustratie van hoe divers en interessant mensen werkelijk zijn. Vandaag nemen we onze eerste duik in de diepe zee van "Philias" om te begrijpen dat we allemaal wel iets geks hebben dat ons opwindt. Ik ben Ashley Hamer, en dit is Taboo Science, de podcast die de vragen beantwoordt die je niet mag stellen.
Definiëren van het 'Abnormale'
Eerst een paar definities. Kink is een overkoepelende term die elke seksuele interesse omvat die buiten de heersende norm valt. Een fetisj is een verhoogde aantrekkingskracht tot bepaalde objecten of lichaamsdelen, met uitzondering van de genitaliën. Volgens deze definities kunnen veel ogenschijnlijk alledaagse seksuele interesses als kinks worden beschouwd. Ben je een billenman? Houd je van katholieke schoolmeisjesuniformen? Raak je extra opgewonden van het dragen van lingerie of leer? Gefeliciteerd! Je bent lid van de kink-club. Je papieren zijn onderweg.
"Je kunt een interesse hebben in iets dat wij kinky noemen," zegt Dr. Christian Joyal. "Bijvoorbeeld vastgebonden worden op je bed. Het wordt als kinky beschouwd, maar voor mij is het gewoon dat je je seksleven wilt opfleuren. Dat is alles. Het is heel goed. Dus dit is gewoon een kinky interesse."
Dr. Christian Joyal is professor aan de psychologieafdeling van de Universiteit van Quebec en fulltime onderzoeker aan het Philippe-Pinel Instituut in Montreal, een forensisch psychiatrisch ziekenhuis. Hij is een grote naam in het veld; meer dan één gast stemde toe om aan deze podcast deel te nemen nadat ik had genoemd dat ik met hem had gesproken.
Een stap verder dan kinky is wat wetenschappers een parafilie noemen. De definitie van parafilie komt uit de DSM, de Diagnostic and Statistical Manual of Mental Disorders. Dit is de grote lijst van psychische stoornissen die niet alleen door professionals in de geestelijke gezondheidszorg wordt gebruikt, maar ook door rechtbanken en wetgevers. Het is dus van groot belang dat deze accuraat is over wat normaal en wat pathologisch is. Specifiek definieert de DSM een parafilie als elke intense en aanhoudende seksuele interesse anders dan seksuele interesse in genitale stimulatie of voorbereidend voorspel met fenotypisch normale, fysiek volwassen en instemmende menselijke partners. De interesse moet gelijk zijn aan of groter zijn dan de interesse in de eerdergenoemde "normale" seks, die ook een naam heeft: normofilie. "Zo noemden ze het vroeger," zegt Dr. Joyal, "maar wat is 'normaal' precies?"
En dan is er nog een stap boven een parafilie. Dr. Joyal legt uit: "Als je daardoor van streek raakt, of als het je dagelijks leven verstoort – je denkt bijvoorbeeld constant aan seks, kijkt porno op je werk – dan heb je een parafiele stoornis."
Het probleem met deze definities is dat ze allemaal afhangen van een solide definitie van 'normaal'. Maar die definitie is gebaseerd op algemene overtuigingen, niet op feitelijke gegevens. Daarom besloten Dr. Joyal en zijn team hun eigen onderzoek te doen. "We begonnen mensen uit de algemene bevolking te interviewen," vertelt hij, "en vroegen hen wat ze wel en niet leuk vonden in hun seksleven." En een heleboel mensen bleken dingen leuk te vinden die de DSM niet als normaal beschouwde.
"Volgens de DSM is een seksuele fantasie voldoende om als parafilie te worden bestempeld," legt Joyal uit. "Maar dat hangt ervan af. Bij het Pinel Instituut werken we met mensen die een diagnose pedofilie hebben of verkrachters zijn. In bepaalde gevallen kan een seksuele fantasie een enorme rode vlag zijn. Maar in de meeste gevallen is dat niet zo." Hij ontdekte bijvoorbeeld dat twee derde van de vrouwen in de algemene bevolking fantasieën heeft over verkrachting. Twee derde! Dit betekent niet dat twee derde van de vrouwen daadwerkelijk verkracht wil worden. Het is een fantasie. "Er is een groot verschil tussen een seksuele fantasie in je hoofd, die bedoeld is om opwinding te wekken, en een seksuele wens of verlangen," benadrukt Joyal. "Vooral vrouwen schreven ons: 'Natuurlijk wil ik dit niet in het echt meemaken.' Maar in hun hoofd raken ze opgewonden van het idee dat hun partner hen met geweld neemt. Zolang je opgewonden raakt van het idee, werkt de fantasie prima en is het perfect. We realiseerden ons dat hoe meer seksuele fantasieën je hebt, hoe meer seksueel tevreden je bent met je seksleven."
Het Onderzoek naar Seksuele Diversiteit
Decennialang heeft de DSM dus gezegd dat dingen als verkrachtingsfantasieën abnormaal zijn, terwijl Dr. Joyal met een simpel onderzoek kan aantonen dat ze fout zitten. Niet alleen zijn deze fantasieën geen schadelijke afwijking, ze verbeteren zelfs het seksleven van mensen. Waar zit de disconnectie? Tot voor kort was er bijna geen financiering voor het soort onderzoek dat Dr. Joyal doet. "Tien jaar geleden was het totaal onmogelijk," zegt hij. "We kregen nooit subsidies voor onderzoek naar seksuele zaken, behalve AIDS, verkrachtingen en dergelijke. Seksuele diversiteit interesseerde niemand." Dus begonnen ze het in hun eigen tijd te doen, in de weekenden, onbetaald, uit pure nieuwsgierigheid.
Maar toen veranderde er iets. "Na de Me Too-beweging veranderde alles," zegt Joyal. "We kregen ineens veel subsidiemogelijkheden om aan seksualiteit in het algemeen te werken."
Dit brengt me bij Aella, een sekswerker en, de laatste tijd steeds meer, een seksuoloog die online onderzoek doet naar fetisjen. Aella is een soort internetberoemdheid met een enorme aanhang. Ik wilde met haar praten omdat ze de auteur is van wat waarschijnlijk de grootste enquête over kinks en fetisjen ooit is, met meer dan 700.000 respondenten in oktober 2023. En ze deed het uit pure nieuwsgierigheid. "Ik heb een paar behoorlijk gestoorde kinks en ik wilde weten hoe gestoord ze zijn," vertelt ze. "Maar om dat te weten, moet je alle andere kinks meten om te kunnen vergelijken. En dat startte een hele reeks vragen." Hoewel ze ze niet publiekelijk deelt uit schaamte, is haar doel juist om die schaamte te verminderen.
Aella heeft geen universitair diploma. Ze leerde statistiek gaandeweg en toen haar dataset te groot werd voor Google Sheets, leerde ze Python. Ondanks haar onconventionele achtergrond gelooft ze dat haar onderzoek van hogere kwaliteit is dan veel van wat er bestaat. "Ik begon onlangs de literatuur door te spitten en dacht: 'Wow, veel van wat ik doe is eigenlijk veel beter dan wat er nu in de literatuur bestaat.' Er is gewoon niet veel. En wat er wel is... het is moeilijk om goede, representatieve steekproeven te krijgen voor zoiets. Mensen willen niet over hun kinks praten. Ik ben misschien niet perfect, maar op dit moment ben ik misschien wel een van de besten die we hebben."
Hoe krijgt Aella mensen zover om haar enquête in te vullen? Ze gebruikte een marketingaanpak. Wanneer je de enquête afrondt, krijg je een "kink-score" en een populair personage dat daarmee overeenkomt, zoals Bambi (minst kinky) of de Geest uit Aladdin (zeer kinky). Dit maakte het leuk en deelbaar. "Ik behandel enquêtes als een product," legt ze uit. "Het resultaat moet je een plezierprikkel geven en je aanzetten om het te delen." Respondenten kregen een personage om te delen, en Aella kreeg een enorme hoeveelheid data.
De Resultaten: Van Tongzoenen tot Incest
Wat maakt iets tot een fetisj? Volgens Aella is het subjectief. "Ik denk dat dingen veel meer op een spectrum liggen. Ik noem alles een 'seksuele interesse'. Hoe taboe en zeldzamer het wordt, hoe groter de kans dat we het als een fetisj zien." De meest populaire en minst taboe seksuele interesse in haar enquête? Tongzoenen. De minst populaire en meest taboe interesse was pedofilie, net iets erger dan necrofilie.
Aella paste ook factoranalyse toe op de data, een statistische methode om een groot aantal variabelen te reduceren tot kleinere groepen. Hiermee identificeerde ze tien clusters van kinks en fetisjen die vaak samen voorkomen:
Machtsdynamiek: Bondage, sm, vernedering, verkrachtingsspel.
Swingen: Voyeurisme, exhibitionisme, anonieme seks, meerdere partners, cuckoldry.
Vanille: Tongzoenen, romantiek, kietelen, sensualiteit.
Monster-liefhebbers: Monsters, bestialiteit, dierlijke transformaties.
Extreme Vrouwelijkheid: Borstimplantaten, hoge hakken.
Geslachtspel: Pegging, sissificatie (waarbij een man een vrouwelijke rol aanneemt).
Extreme Mannelijkheid: Homoseksualiteit, onderdanigheid, pijn ontvangen.
Horror: Bloed, zombies, necrofilie, brutaliteit.
Afkeer: Scat (ontlasting), urine, luiers, oksels, scheten, voeten.
Incest en Leeftijdsspel: Waar mensen de rol van een andere leeftijdsgroep aannemen.
Het is belangrijk te onthouden: fantasieën zijn geen realiteit.
De Oorsprong van Kink
Waarom raken mensen opgewonden van dingen die we evolutionair gezien zouden moeten vermijden, zoals pijn, enge dingen of vernedering? Volgens Dr. Joyal beginnen fundamentele seksuele interesses vaak in de kindertijd. "Mensen die bijvoorbeeld exclusief homoseksueel zijn, wisten dit vaak al op achtjarige leeftijd," zegt hij. "Hetzelfde geldt voor mensen met exclusieve pedofilie; als kind voelden ze zich aangetrokken tot leeftijdsgenoten, en die aantrekking veranderde niet naarmate ze ouder werden." Hij merkt op dat veel mensen die van BDSM houden, zich herinneren dat ze als kind al gefascineerd waren door bijvoorbeeld bondage in Kuifje-strips. "Kuifje wordt in elk boek vastgebonden," zegt Joyal. "Veel van onze respondenten vertelden ons dat ze hierdoor opgewonden raakten, nog niet seksueel, maar ze vonden het wel heel interessant."
Het is niet zo dat Kuifje deze mensen in BDSM geïnteresseerd maakte. De interesse was er waarschijnlijk al en vond toevallig een uitlaatklep in deze stripboeken. Aella waarschuwt voor het verkeerd toeschrijven van oorzaken. "Het is heel gemakkelijk om de oorzaak verkeerd te interpreteren," zegt ze. "Misschien is het genetisch. Misschien zijn voetfetisjen om de een of andere reden aangeboren." We weten het simpelweg niet zeker vanwege het gebrek aan onderzoek.
Dr. Joyal voegt toe: "Fundamentele interesses beginnen meestal in de kindertijd. Daarna ontdek je secundaire voorkeuren tijdens je studententijd, en een derde soort voorkeur ontdek je met een partner als je bijvoorbeeld 30 bent."
Trauma, Schaamte en Heling
Dit zou goed nieuws moeten zijn voor iedereen die zich schaamt voor zijn kinks. Je bent niet beschadigd; je hebt gewoon een seksuele interesse die buiten de 'norm' valt. Maar veel mensen zijn als kind misbruikt, en sommigen hebben kinks die dat misbruik weerspiegelen. En dat is oké. Aella stelt: "Zelfs als een fetisj wordt veroorzaakt door misbruik in de kindertijd, voelt dat voor mij niet schaamtevol. Het is een volkomen rationele reactie van je lichaam om te proberen je te beschermen."
Voor sommigen kan het naspelen van deze scenario's zelfs een manier zijn om te helen. "Voor sommigen is het echt geweldig," zegt Dr. Joyal, "omdat ze ons vertellen dat ze een traumatische situatie naspelen, maar nu hebben ze de controle over die situatie. Het is geruststellend en goed voor hen."
De schadelijke boodschappen over kink komen niet alleen uit de media, maar ook van wetenschappelijke instanties. De definitie van een parafiele stoornis is circulair: als het leed veroorzaakt, is het een stoornis. Maar de diagnose zelf kan dat leed juist veroorzaken. "Precies!" beaamt Dr. Joyal. "Wat je zegt is zo waar. Het is zo circulair. We maken mensen van streek die niet van streek zijn. Het is een grap. Over 25 jaar zullen we zeggen: 'We waren zo dom in 2023 dat we probeerden te definiëren wat normaal en abnormaal is in de seksuologie.' Het zal hetzelfde zijn als met homoseksualiteit in de DSM-II."
Aella wil er zeker van zijn dat niemand zich van streek voelt vanwege zijn kink. "Er is niet veel empathie," zegt ze. "Ik wil meer compassie en begrip brengen. Als je een vreemde fetisj hebt en meer over jezelf wilt weten, is er doodse stilte. Niemand weet iets. En dat is gewoon triest. Ik wil de mogelijkheid bieden om zonder schaamte naar jezelf te kijken en te zeggen: 'Cool. Hier is informatie. Iemand gaf genoeg om jou om dit te doen.'"
En dat is wat ik hoop te doen met Philias. We gaan de kinks die je misschien kent als raar, schaamtevol of als de clou van een grap, in het licht houden. Als je je met deze kinks identificeert, leer je meer over jezelf. En als dat niet zo is, ontdek je hoeveel we allemaal gemeen hebben. Iedereen is gewoon op zoek naar een beetje vreugde in zijn leven.
